% Source: https://artofproblemsolving.com/community/c3h1388504_asterisks_equation
% Stable identity: aops-topic-1388504-q01
Let $T$ be the sum of the two tens digits on the left and $U$ the sum of the other six digits on the left, and let the right side be $\overline{ij}=10i+j$. Then the equation is $10T+U=10i+j$, and since all ten digits are used, $T+U+i+j=45$.

\textbf{Digit sum of the right side.} Modulo $9$, $10T+U\equiv T+U=45-(i+j)\equiv-(i+j)$, while $10i+j\equiv i+j$. So $2(i+j)\equiv0$, hence $9\mid i+j$. As $i\ne j$ are digits with $i\ge1$, $1\le i+j\le17$, so $i+j=9$.

\textbf{The tens digit.} Now $T+U=36$, so the left side is $10T+U=9T+36$, and the right side is $10i+(9-i)=9i+9$. Thus $i=T+3$. The two tens digits on the left are different and nonzero, so $T\ge1+2=3$ and $i\ge6$. So the right side is one of $63$, $72$, $81$, $90$.

\textbf{All four occur:}
\[
10+24+5+7+8+9=63,\quad 10+34+5+6+8+9=72,\quad 20+34+5+6+7+9=81,\quad 12+53+4+6+7+8=90 .
\]
So the possible right-hand sides are $\boxed{63,\ 72,\ 81,\ 90}$.
